\subsubsection{The spectrum frame from radical ideals}
\label{sec:p10-radical-ideal-spectrum}

A construction of opens need not begin with an object of points.
The following construction uses ideals directly and works in the
internal intuitionistic language of the topos.
An ideal is an additive subgroup closed under multiplication by
arbitrary ring elements. Call an ideal \(I\) radical when
\(a^2\in I\) implies \(a\in I\).
For a subobject \(U\hookrightarrow A\), write \(\rho(U)\) for the
intersection of all radical ideals containing \(U\).
This intersection exists using the power object and universal
quantification already constructed; the family is inhabited by \(A\).
For a ring element \(a\), put \(D(a)=\rho(\{a\})\).

\begin{hbtheorem}{thm:p10-radical-spectrum-frame}{proved=annex:p10-radical-spectrum-frame}{The frame of radical ideals}
The radical ideals of \(A\), ordered by inclusion, form an internal
frame. Meets are intersections and the join of a family is the radical
closure of its union. The bottom element is \(\rho(\varnothing)\)
and the top is \(A\).
Its principal generators satisfy
\[
 \begin{aligned}
 D(0)&=0,\qquad D(1)=1,\\
 D(ab)&=D(a)\wedge D(b),\\
 D(a+b)&\leq D(a)\vee D(b).
 \end{aligned}
\]
Here \(0,1\) on the right are frame bounds, not ring elements.
Every radical ideal \(I\) is the join of the \(D(a)\) with \(a\in I\).
\end{hbtheorem}

\begin{hbtheorem}{thm:p10-radical-spectrum-presentation}{proved=annex:p10-radical-spectrum-presentation}{Universal property of the generators}
Let \(H\) be any internal frame, and let \(d:A\to H\) satisfy
\[
 d(0)=0,\quad d(1)=1,\quad
 d(ab)=d(a)\wedge d(b),\quad
 d(a+b)\leq d(a)\vee d(b).
\]
There is exactly one frame homomorphism from the radical-ideal frame
to \(H\) that sends \(D(a)\) to \(d(a)\). Its value at \(I\) is
\[
 I\longmapsto\bigvee_{a\in I}d(a).
\]
A frame homomorphism preserves arbitrary joins and finite meets,
including the top element.
\end{hbtheorem}

These equations describe the points of the frame.
The frame is constructed without assuming that points exist.
Taking \(H\) to be the frame of truth values identifies its
frame homomorphisms with the prime predicates of
\cref{def:p10-lawvere-sheaves-prime}.
An absence of such predicates therefore does not collapse the
construction into the topology of an empty object.
The proof makes no appeal to the Boolean prime ideal theorem.
